%\documentclass[11pt,a4paper,uplatex,dvipdfmx]{ujarticle} 		% for uplatex
\documentclass[11pt,a4j,dvipdfmx]{jarticle} 					% for platex
\input{pieces/form00_header} % pieces
\input{pieces/kakenhi7} % pieces
\input{pieces/form01_dcpd_header} % pieces
\input{pieces/hook3} % pieces
%#Name: orf
\input{pieces/form03_dcpd_headers} % pieces
\input{pieces/form04_orf_header} % pieces
% ===== Global definitions for the Kakenhi form ======================
%　基本情報
%
%------ 研究課題名  -------------------------------------------
\newcommand{\研究課題名}{象の卵}

%----- 研究機関名と研究代表者の氏名-----------------------
\newcommand{\研究機関名}{逢坂大学}
\newcommand{\研究代表者氏名}{湯川秀樹}
\newcommand{\me}{\underline{\underline{H.~Yukawa}}} 
\input{pieces/inst_dcpd} % pieces
% user07_header
% ===== my favorite packages ====================================
% ここに、自分の使いたいパッケージを宣言して下さい。
\usepackage{wrapfig}
%\usepackage{amssymb}
%\usepackage{mb}
%\DeclareGraphicsRule{.tif}{png}{.png}{`convert #1 `dirname #1`/`basename #1 .tif`.png}
\usepackage{lineno}

% ===== my personal definitions ==================================
% ここに、自分のよく使う記号などを定義して下さい。
\newcommand{\klpionn}{K_L \to \pi^0 \nu \overline{\nu}}
\newcommand{\kppipnn}{K^+ \to \pi^+ \nu \overline{\nu}}

% ----- 業績リスト用 -------------
\newcommand{\paper}[6]{%
	% paper{title}{authors}{journal}{vol}{pages}{year}
	\item ``#1'', #2, #3 {\bf #4}, #5 (#6).			% お好みに合わせて変えてください。
}

\newcommand{\etal}{\textit{et al.\ }}
\newcommand{\ca}[1]{*#1}	% corresponding author;   \ca{\yukawa}  みたいにして使う
\newcommand{\invitedtalk}{招待講演}

\newcommand{\yukawa}{H.~Yukawa}					% no underline
%\newcommand{\yukawa}{\underline{\underline{H.~Yukawa}}}	% with 2 underlines
\newcommand{\tomonaga}{S.~Tomonaga}

\newcommand{\prl}{Phys.\ Rev.\ Lett.\ }		% よく使う雑誌も定義すると楽

% ===== 欄外メモ ==================
\newcommand{\memo}[1]{\marginpar{#1}}
%\renewcommand{\memo}[1]{}	% 全てのメモを表示させないようにするには、行頭の"%"を消す


%\input{../../sample/simple/contents}	% skip
\input{pieces/hook5} % pieces

\begin{document}
\input{pieces/hook7} % pieces
%#Split: 01_background  
%#PieceName: p01_background
\input{pieces/p01_background_00}
\section{研究の位置づけ}
%　　　　＜＜最大　1ページ＞＞

%s03_background
%begin 本研究の着想に至った経緯など ＝＝＝＝＝＝＝＝＝＝＝＝＝＝＝＝＝＝＝＝
		風呂で巨大な温泉卵について考えていて、ふと思いついた。

	準備はしようとしている。

	唯一無二。
%end 本研究の着想に至った経緯など ＝＝＝＝＝＝＝＝＝＝＝＝＝＝＝＝＝＝＝＝

\input{pieces/p01_background_01}

%#Split: 02_purpose_plan  
%#PieceName: p02_purpose_plan
\input{pieces/p02_purpose_plan_00}
\section{研究目的・内容等}
%　　　　＜＜最大　2ページ＞＞

%s02_purpose_plan_dcpd
%begin 研究目的と研究計画short留意事項なし ＝＝＝＝＝＝＝＝＝＝＝＝＝＝＝＝＝＝＝＝
\textbf{\\　　　　　＊＊＊　以下は、あくまで例です。真似しないでください。　＊＊＊\\
　　　　　＊＊＊　本文はもちろん、節の切り方や論理の組み方は　　　＊＊＊\\
　　　　　＊＊＊　ご自分の気に入ったスタイルで書いてください。　　＊＊＊}

	 象の卵の研究目的は．．．

	象の卵の研究計画は．．．

	\vspace{1cm}
	\begin{thebibliography}{99}
		\bibitem{teramura} 寺村輝夫、「ぼくは王様 - ぞうのたまごのたまごやき」.
	\end{thebibliography}
%end 研究目的と研究計画short留意事項なし ＝＝＝＝＝＝＝＝＝＝＝＝＝＝＝＝＝＝＝＝
\input{pieces/p02_purpose_plan_01}

%#Split: 03_why_abroad  
%#PieceName: p03_why_abroad
\input{pieces/p03_why_abroad_00}
\section{外国で研究することの意義}
%　　　　＜＜最大　1ページ＞＞

% s07_why_abroad
%begin 外国で研究する事の意義 ＝＝＝＝＝＝＝＝＝＝＝＝＝＝＝＝＝＝＝＝
	私は今まで、象の卵の可能性について主に文献を漁って研究をしてきた。
	そうした長年の研究の末分かったことの一つは、日本に現在、自然界に生息
	する象はいないということである。
	最も最近生息した象はケナガマンモスのようであるが、
	祖父が子供の頃には既に絶滅していたそうである。
	マンモスの氷漬けの個体は北海道で見つかったが、卵は見つかっていない。
	また最近では2005年に愛知県のある会場で氷漬けの個体が見つかったが、
	これは実は密かにロシアから持ち込まれたものであり、国産象ではない。
	
	こうした経験から、象の卵を日本で探していても見つからないということを
	強く実感し、海外で研究する決心をした次第である。
	特に、象の卵を探す夢を子供の頃に私に与えてくれた
	Dr. Seussにぜひとも指導を仰ぎたく、師の元に行って研究を行う。
%end 外国で研究する事の意義 ＝＝＝＝＝＝＝＝＝＝＝＝＝＝＝＝＝＝＝＝

\input{pieces/p03_why_abroad_01}

%#Split: 04_rights  
%#PieceName: p04_rights
\input{pieces/p04_rights_00}
\section{人権の保護及び法令等の遵守への対応}
%　　　　＜＜最大　1ページ＞＞

% s09_rights
%begin 人権の保護及び法令等の遵守への対応 ＝＝＝＝＝＝＝＝＝＝＝＝＝＝＝＝＝＝＝＝
	象の卵のES細胞の培養、象のクローンの生成などは行わない。

	\LaTeX の便利な機能については、\texttt{egg\_***.tex} や\texttt{sample\_pdf/egg\_***.pdf}をご覧ください。
%end 人権の保護及び法令等の遵守への対応 ＝＝＝＝＝＝＝＝＝＝＝＝＝＝＝＝＝＝＝＝

\input{pieces/p04_rights_01}

%#Split: 05_abilities  
%#PieceName: p05_abilities
\input{pieces/p05_abilities_00}
\section{研究遂行力の自己分析}
%　　　　＜＜最大　2ページ＞＞

% s14_abilities
%begin 自己分析 ＝＝＝＝＝＝＝＝＝＝＝＝＝＝＝＝＝＝＝＝
\DCPDInstructionsA\\% <-- 留意事項：これは消すか、コメントアウトしてください。
\noindent
\textbf{(1) 研究に関する自身の強み}
\DCPDInstructionsB% <-- 留意事項：これは消すか、コメントアウトしてください。

	応募者は過去20年間、７つの海を隅から隅まで航海し、
	浅瀬から深海まで潜り、文字通り東西南北上下の３次元で
	シロナガスクジラの卵の探索を行ってきた(業績\ref{pub:whale})。
	シロナガスクジラに飲み込まれそうになったり、海賊に捕まるなどの危険な目にも
	あったが、それにもめげず、研究を遂行してきた強靭な能力を有する。
	\begin{enumerate}
		\paper{Search for whale eggs}{\yukawa\ \etal}{Rev.\ Oceanic Mysteries}{888}{99}{2017}
			\label{pub:whale}
				
		\paper{Theory of Elephant Eggs}{\yukawa, Kara Juzo \etal}{\prl}{800}{800-804}{2005}
			\label{pub:theoegg}
				
		\paper{仔象は死んだ}{Kobo Abe}{安部公房全集}{26}{100-200}{2004}
		
		\paper{The Elephant's Child (象の鼻はなぜ長い)}{R.~Kipling}{Nature}{999}{777-799}{2003}

		\paper{You can't Lay an Egg If You're an Elephant}{F.~Ehrlich}
			{JofUR\\({\tt www.universalrejection.org})}{{\bf N/A}}{2002}
		
		% 下のように書いてもいいけど、めんどくさいし、表示の仕方を変えようとしたら大変。
		\item ``Egg of Elephant-Bird'', 
				\underline{A.~Cooper},
				Nature, {\bf 409}, 704-707 (2001).	% 	
		\input{jack_pub}	% << only for demonstration. Please delete it or comment it out.
	\end{enumerate}
\vspace{5mm}
\noindent
\textbf{(2) 今後研究者として更なる発展のため必要と考えている要素}

研究費を獲得する術。

%end 自己分析 ＝＝＝＝＝＝＝＝＝＝＝＝＝＝＝＝＝＝＝＝

\input{pieces/p05_abilities_01}

%#Split: 99_tail
\input{pieces/hook9} % pieces
\end{document}

